\section{引言}

\subsection{编写目的}
本文档用于描述学术AI鉴伪系统在迭代开发阶段的概要设计方案，明确系统的总体结构、模块划分、关键数据关系、主要交互流程以及部署与质量设计，为后续详细设计、编码实现、测试验证和系统演进提供依据。

\subsection{文档范围}
本文档重点覆盖以下内容：
\begin{itemize}[leftmargin=2em]
    \item 迭代后系统的总体架构与模块边界；
    \item 用户端、管理端、后端服务与AI推理服务之间的协作关系；
    \item 学术图像检测、全篇论文检测、Review检测、人工审核与报告生成的设计方案；
    \item 核心数据实体、接口分层、关键业务流程与运行部署设计；
    \item 可维护性、可扩展性、安全性与性能等质量属性设计。
\end{itemize}

\subsection{参考资料}
\begin{itemize}[leftmargin=2em]
    \item 《学术AI鉴伪系统需求规格说明书（迭代版）》；
    \item 当前项目代码仓库中的前端、后端与AI服务实现；
    \item 项目现有接口文档、实体表与交接说明。
\end{itemize}

\subsection{术语说明}
\begin{longtable}{|p{4cm}|p{10cm}|}
\hline
术语 & 说明 \\ \hline
学术资源 & 系统支持处理的业务对象，包括学术图像、论文文件与同行评审Review。 \\ \hline
检测任务 & 用户针对一个或多个学术资源发起的一次检测处理请求。 \\ \hline
检测结果 & 系统对指定资源执行自动检测后产生的结构化输出。 \\ \hline
人工审核 & 在自动检测基础上，由人工参与复核、判断或补充说明的业务过程。 \\ \hline
综合报告 & 汇总自动检测与人工审核信息后形成的可下载输出结果。 \\ \hline
\end{longtable}
